% (c) Nikita Lisitsa, lisyarus@gmail.com, 2026

\documentclass[10pt]{beamer}

\usepackage[T2A]{fontenc}
\usepackage[russian]{babel}
\usepackage{minted}

\usepackage[most]{tcolorbox}
\tcbuselibrary{minted}

\usepackage{graphicx}
\graphicspath{ {./images/} }

\usepackage{adjustbox}

\usepackage{color}
\usepackage{soul}
\usepackage{multicol}
\usepackage{hyperref}
\usepackage{tikz}
\usetikzlibrary{arrows.meta,positioning,calc}

\usetheme{metropolis}

\definecolor{red}{rgb}{1,0,0}
\definecolor{codebg}{RGB}{29,35,49}
\setminted{fontsize=\footnotesize}

\makeatletter
\newcommand{\slideimage}[1]{
  \begin{figure}
    \begin{adjustbox}{width=\textwidth, totalheight=\textheight-2\baselineskip-2\baselineskip,keepaspectratio}
      \includegraphics{#1}
    \end{adjustbox}
  \end{figure}
}
\makeatother

\title{Язык программирования C++}
\subtitle{Введение в курс}
\author{Лисица Никита}
\date{2026}

\setbeamertemplate{footline}[frame number]

\usemintedstyle{tango}

\begin{document}

\frame{\titlepage}

\begin{frame}[fragile]
\frametitle{Цели курса}
\begin{itemize}
\item Узнать, \textbf{как работает компьютер}: регистры, стек, счётчик команд
\pause
\item Изучить подробно \textbf{языки C и C++}
\pause
\item Рассмотреть \textbf{парадигмы программирования} на примере C и C++: процедурное, объектно-ориентированное, обобщённое, функциональное
\pause
\item Рассмотреть \textbf{шаблоны проектировния} на примере C и C++
\end{itemize}
\end{frame}

\begin{frame}[fragile]
\frametitle{Пререквезиты курса}
\begin{itemize}
\item \textbf{Желательно:} умение программировать на каком-нибудь языке (Basic, Pascal, Python, ...)
\pause
\item \textbf{Обязательно:} готовность разобраться с базовым синтаксисом C (переменные, условия, циклы, функции)
\end{itemize}
\end{frame}

\begin{frame}[fragile]
\frametitle{Зачем изучать C и C++?}
\begin{itemize}
\item \textbf{Низкоуровневые языки} \(\Longrightarrow\) помогают лучше понять, как работает `железо'
\pause
\item \textbf{Высокопроизводительные языки} \(\Longrightarrow\) на них нишут операционные системы, драйверы, базы данных, текстовые редакторы, веб-серверы, мультимедиа-редакторы, видеоигры
\pause
\item \textbf{Старые и популярные языки} \(\Longrightarrow\) синтаксис многих современных языков основан на C
\end{itemize}
\end{frame}

\begin{frame}[fragile]
\frametitle{План курса}
\begin{itemize}
\item \textbf{Первый семестр}: мини-курс (8 занятий) по языку C
\pause
\item \textbf{Второй семестр}: полный курс (15 занятий) по языку C++
\end{itemize}
\end{frame}

\begin{frame}[fragile]
\frametitle{Отчётность по курсу}
\begin{itemize}
\item \textbf{Лабораторные работы}: маленькие задания, выдаются \textbf{на каждой практике}
\item \textbf{Домашние работы}: большие задания, выдаются \textbf{на практиках} 1-2 раза в семестр
\item \textbf{Письменные тесты}: выдаются \textbf{на лекциях} 2-3 раза в семестр
\item \textbf{Экзамен}: 1 раз в конце семестра
\item За все виды отчётности начисляются баллы
\item Для допуска на экзамен нужно
\begin{itemize}
\item Набрать \(\geq30\%\) общего количества баллов без экзамена
\item Набрать \(>0\) баллов за домашние работы
\end{itemize}
\item Для сдачи практических заданий нужно зарегистрироваться в системе \nolinkurl{gitlab}
\end{itemize}
\end{frame}

\begin{frame}[fragile]
\frametitle{Рекомендуемые инструменты}
\begin{itemize}
\item Linux (напр. Ubuntu)
\pause
\item Терминал
\pause
\item Компилятор \nolinkurl{gcc}, утилита \nolinkurl{make}
\pause
\item \textbf{N.B.:} Если вы понимаете, что делаете, можно и другое ПО, но проверяться будет именно на этом
\end{itemize}
\end{frame}

\end{document}
